\documentclass[11pt,a4paper]{article}
\usepackage[utf8]{inputenc}
\usepackage[margin=1in]{geometry}
\usepackage{hyperref}
\usepackage{graphicx}
\usepackage{amsmath}
\usepackage{booktabs}

\title{RSCT Review: PONTE: Personalized Orchestration for Natural Language Trustworthy Explanations}
\author{Swarm-It Research Discovery\\\small Automated RSCT-Based Analysis}
\date{March 09, 2026}

\begin{document}
\maketitle

\begin{abstract}
This document provides an automated review of the paper ``PONTE: Personalized Orchestration for Natural Language Trustworthy Explanations'' from an RSCT (Representation-Solver Compatibility Testing) perspective. The paper achieved an RSCT relevance score of 34.3% and topic similarity of 44.1%.
\end{abstract}

\section{Paper Information}
\begin{itemize}
\item \textbf{Title:} PONTE: Personalized Orchestration for Natural Language Trustworthy Explanations
\item \textbf{Authors:} Vittoria Vineis, Matteo Silvestri, Lorenzo Antonelli, Filippo Betello, Gabriele Tolomei
\item \textbf{Source:} \url{https://arxiv.org/abs/2603.06485v1}
\item \textbf{RSCT Relevance:} 34.3%
\item \textbf{Topic Match:} 44.1%
\item \textbf{Key Concepts:} hallucination, alignment
\end{itemize}

\section{Original Abstract}
Explainable Artificial Intelligence (XAI) seeks to enhance the transparency and accountability of machine learning systems, yet most methods follow a one-size-fits-all paradigm that neglects user differences in expertise, goals, and cognitive needs. Although Large Language Models can translate technical explanations into natural language, they introduce challenges related to faithfulness and hallucinations. To address these challenges, we present PONTE (Personalized Orchestration for Natural language Trustworthy Explanations), a human-in-the-loop framework for adaptive and reliable XAI narratives. PONTE models personalization as a closed-loop validation and adaptation process rather than prompt engineering. It combines: (i) a low-dimensional preference model capturing stylistic requirements; (ii) a preference-conditioned generator grounded in structured XAI artifacts; and (iii) verification modules enforcing numerical faithfulness, informational completeness, and stylistic alignment, optionally supported by retrieval-grounded argumentation. User feedback iteratively updates the preference state, enabling quick personalization. Automatic and human evaluations across healthcare and finance domains show that the verification-refinement loop substantially improves completeness and stylistic alignment over validation-free generation. Human studies further confirm strong agreement between intended preference vectors and perceived style, robustness to generation stochasticity, and consistently positive quality assessments.

\section{RSCT Analysis}
This paper presents research with 34% relevance to RSCT theory.

\textbf{Abstract Summary:} Explainable Artificial Intelligence (XAI) seeks to enhance the transparency and accountability of machine learning systems, yet most methods follow a one-size-fits-all paradigm that neglects user differences in expertise, goals, and cognitive needs. Although Large Language Models can translate technical explanations into natural language, they introduce challenges related to faithfulness and hallucinations. To address these challenges, we present PONTE (Personalized Orchestration for Natural lan...

\textbf{RSCT Relevance:} The work touches on concepts related to hallucination, alignment, which have connections to the Representation-Solver Compatibility Testing framework.

\textbf{Further Analysis:} A detailed manual review is recommended to fully assess the implications for RSCT-based certification approaches.


\subsection*{RSCT Certification Metrics}
\begin{tabular}{ll}
\textbf{Relevance (R):} & 0.375 \\
\textbf{Spurious (S):} & 0.318 \\
\textbf{Noise (N):} & 0.307 \\
\textbf{Kappa ($\kappa$):} & 0.550 \\
\end{tabular}

The kappa score of 0.550 indicates moderate representation-solver compatibility.


\section{Relevance to Swarm-It}
This paper was identified by the Swarm-It research discovery pipeline as potentially relevant to RSCT-based AI certification. The combined score of 38.2% places it in the moderate relevance category for further investigation.

\vspace{1em}
\hrule
\vspace{0.5em}
\small{Generated by Swarm-It Research Discovery | \url{https://swarms.network} | RSCT Certified}

\end{document}
